\documentclass[10pt,a4paper]{amsart}
\usepackage[margin=19mm]{geometry}
\usepackage{amsmath,amssymb,mathtools}
\usepackage[scr=rsfs,cal=euler]{mathalfa}
\usepackage{xfrac}
\usepackage[unicode,hidelinks,bookmarks=false]{hyperref}
\newcommand{\ie}{i.e.\ }
\newcommand{\cf}{cf.\ }
\newcommand{\resp}{respectively\ }
\newcommand{\Subsection}{\subsection}
\providecommand{\nobreakdash}{\nobreak\hbox{-}}
\setlength{\emergencystretch}{2em}
\multlinegap0pt
\usepackage{amssymb}
\usepackage{xcolor}
\usepackage{bm}
\usepackage[shortlabels]{enumitem}
\usepackage{mathtools}
\usepackage{tikz}
\usepackage{float}
\newtheorem{assumption}{Assumption}
\newtheorem{lemma}{Lemma}
\def\C{\mathbb{C}}
\def\E{\mathbb{E}}
\def\P{\mathbb{P}}
\def\R{\mathbb{R}}
\def\Tr{\operatorname{Tr}}
\def\bD{\mathbf{D}}
\def\cA{\mathcal{A}}
\def\cC{\mathcal{C}}
\def\cL{\mathcal{L}}
\def\cP{\mathcal{P}}
\def\cY{\mathcal{Y}}
\def\cZ{\mathcal{Z}}
\def\de{\operatorname{d}}
\def\eps{\varepsilon}
\def\op{\mathrm{op}}
\def\u{\mathbf{u}}
\def\x{\mathbf{x}}
\title{The anisotropic local law for sample covariance matrices\\ under quadratic-form concentration}
\author{Renyuan Ma, Theodor Misiakiewicz$^{\dagger}$}
\date{Source: arXiv:2609.09440v2 (CC BY 4.0)}
\begin{document}
\maketitle

\noindent\emph{Source and licence.} The anisotropic local law for sample covariance matrices\\ under quadratic-form concentration. Version arXiv:2609.09440v2 (CC BY 4.0), distributed by arXiv under the Creative Commons Attribution 4.0 International licence, http://creativecommons.org/licenses/by/4.0/, as recorded in the arXiv licence metadata for this exact version and shown on the abstract page https://arxiv.org/abs/2609.09440v2. Full licence notice: \texttt{licenses/arxiv-2609.09440-CC-BY-4.0.txt}.\par\smallskip

\noindent\textit{Editorial scope.} Counted unit: the article's lemma lem:L\_i-bounds (source0.tex lines 1421--1438). The article's complete original proof of it is delivered with it (source0.tex lines 1439--1551, one \texttt{proof} environment of 113 lines). No mathematics is written, repaired or re-derived: the statement and the proof are the article's own lines, in the article's own notation, and no retained line is substituted or re-flowed.  Retained uncounted as the source-local context the proof needs: lines 314--316; lines 318--327; lines 1279--1297; lines 1317--1329.  Nothing was omitted from any retained span, so the package carries no omission record.  No retained line contains a citation, so the wrapper carries no bibliography and no bibliography entry.  No retained line contains a figure environment, an image-inclusion command or drawing code, so no image is delivered and the package declares no asset dependency.  Editorial formatting only, in addition: an AMS \texttt{amsart} preamble that restates the article's own declarations and macros used by the retained lines, namely the environments \texttt{assumption}, \texttt{lemma} on separate counters, together with the standard packages those lines need.  The retained environments print in this document as Assumption 1, Lemma 1 in order of appearance, whereas the article prints the same items with the numbering of its own full document.  The complete native source of the article as distributed in the arXiv e-print for this exact version is bundled unchanged as \texttt{sources/209814/source0.tex}.
\begin{assumption}[Basic assumptions]\label{ass:A1}
    There exist constants $c,C >0$ such that $c\le n/N\le C$, $\| \Sigma \|_\op \leq C$, and $\Sigma$ is positive definite with at most $(1-c)n$ of its eigenvalues belonging to $[0,c]$.
\end{assumption}

\begin{assumption}[Quadratic-form concentration]\label{ass:A2}
 The vectors $\x_1,\ldots,\x_N\in\R^n$ are independent with $\E\x_i = 0$ and $\E\x_i \x_i^* = \Sigma$. For each $k \geq 1$, there exists a constant $C_k>0$ such that, uniformly in $i$,
\[
 \E\|\x_i\|_2^k\le n^{C_k}.
\]
Finally, for any $\eps,D >0$, there exists a constant $C \equiv C(\eps,D)>0$ such that, for every $i = 1, \ldots, N$ and every $A\in\C^{n\times n}$,
\begin{equation}\label{eq:assumption-qf}
 \P \big[ |\x_i^*A\x_i-\Tr(\Sigma A)| > N^\eps \|A\|_F \big] \le C N^{-D}.
\end{equation}
\end{assumption}

\begin{lemma}[OU-generator]\label{lem:OU-generator}
For $t> 0$, define
\[
    X_t:=e^{-t/2}X_0+\sqrt{1-e^{-t}}G\in\R^{n\times N},\quad X_t=\begin{bmatrix}
        | & \cdots &|\\
        \x_1 & \cdots & \x_N\\
        | & \cdots & |
    \end{bmatrix}.
\]
Here, the columns of $G$ are independent Gaussian vectors with mean zero and covariance $\Sigma$, and $X_0$ is independent of $G$. Let $f:\R^{n\times N}\rightarrow\C$ be a twice continuously differentiable function whose first two derivatives have at most polynomial growth. Then
\[
    \frac{\de }{\de t}\E f(X_t)=\sum_{i=1}^N\E\cL_if(X_t),
\]
where
\[
    \cL_i:=\frac{1}{2}\sum_{\alpha,\beta=1}^n\Sigma[\alpha,\beta]\frac{\de^2}{\de\x_i[\alpha]\de\x_i[\beta]}-\frac{1}{2}\sum_{\alpha=1}^n\x_i[\alpha]\frac{\de}{\de\x_i[\alpha]}
\]
is the OU-generator for the $i$-th column.
\end{lemma}

\begin{lemma}\label{lem:product-and-chain-rule}
Under the hypotheses of Lemma \ref{lem:OU-generator}, let $f,g:\R^{n}\rightarrow\C$ and $h:\C\rightarrow\C$ be such that $h$ is holomorphic on an open neighborhood of the range of $f$, and $f,g$, and $h(f)$ are twice differentiable functions whose first two derivatives have at most polynomial growth. Then, for all $i\in[N]$, the following hold:
\begin{enumerate}
    \item (Product rule)
    \[
        \cL_i(fg)=f\cL_ig+g\cL_if+\nabla f^\top\Sigma\nabla g.
    \]
    \item (Chain rule)
    \[
        \cL_ih(f)=h'(f)\cL_if+\frac{1}{2}h''(f)\nabla f^\top\Sigma\nabla f.
    \]
\end{enumerate}
\end{lemma}

\begin{lemma}\label{lem:L_i-bounds}
Under Assumptions \ref{ass:A1} and \ref{ass:A2}, let $\bD$ be a regular domain and fix $h,t\geq 0$. Suppose there exist $\delta>0$ and a deterministic $\Phi:\bD\rightarrow[N^{-10},1]$ such that, uniformly over all $i\in[N]$, all $z\in\bD$, and all deterministic unit vectors $\u\in\R^{n+N}$, we have
\[
    N^{-1/2}\|B_i\u\|_2, N^{-1}\|A_i\|_F\prec\Phi,\quad \|\Sigma^{1/2}B_i\u\|_2, N^{-1/2}\|\Sigma^{1/2}A_i\|_F\prec N^\delta\sqrt{\Lambda_h},
\]
and
\[
    |\beta+N^{-1}\Tr\Sigma A_i|^{-1}, |\beta+N^{-1}\x_i^*A_i\x_i|^{-1}\prec N^\delta.
\]
Then, for every fixed $L\geq 1$, uniformly over all $l_1,l_2,l_3,l_4\geq 0,l_1+l_2+l_3+l_4=L$, all $i\in[N]$, all $z\in\bD$, and all deterministic unit vectors $\u\in\R^{N+n}$,
\[
    |\E_i\cL_i[\cP_i^L(\u[i]\cY_i)^{l_1}\cZ_i^{l_2}(\u[i]^2\cA_i)^{l_3}(\cC_i\cA_i)^{l_4}]|\prec c_i(\Lambda_h+1) N^{\delta (L+4)}\Phi^{L}.
\]
Here
\[
    |c_i|\leq N^{-1}+N^{-1/2}|\u[i]|+\u[i]^2.
\]
\end{lemma}

\begin{proof}
We first record how the operators $\cL_i$ and $\nabla$ act on the factors $\cY_i,\cZ_i,\cA_i$. This will streamline the later combinatorial arguments. We have
\[
    \cL_i\cY_i=-\frac{1}{2}\cY_i,\quad \cL_i\cA_i=\cL_i\x_i^*(N^{-1}A_i)\x_i=-\cA_i,\quad \cL_i\cZ_i=-\cZ_i
\]
\[
    \nabla\cY_i=N^{-1/2}B_i\u,\quad \nabla\cZ_i=2\cY_i(N^{-1/2}B_i\u),\quad \nabla\cA_i=\nabla\x_i^*(N^{-1}A_i)\x_i=2(N^{-1}A_i\x_i)
\]
\[
    \cL_i\cP_i^L=L\cP_i^{L+1}\cA_i+2L(L+1)\cP_i^{L+2}\x_i^*(N^{-1}A_i)\Sigma(N^{-1}A_i)\x_i,\quad \nabla\cP_i^L=-2L\cP_i^{L+1}(N^{-1}A_i\x_i).
\]
These identities follow directly from the definition of $\cL_i$. We will also repeatedly use the product and chain rules from Lemma \ref{lem:product-and-chain-rule}.

By the product and chain rules, we have
\begin{align*}
    &|\E_i\cL_i[\cP_i^L(\u[i]\cY_i)^{l_1}\cZ_i^{l_2}(\u[i]^2\cA_i)^{l_3}(\cC_i\cA_i)^{l_4}]|=\underbrace{(|\u[i]|)^{l_1+2l_3}|\cC_i|^{l_4}}_{\omega_i}|\E_i\cL_i[\cP_i^L\cY_i^{l_1}\cZ_i^{l_2}\cA_i^{l_3+l_4}]|\\
    &\leq\omega_i|\E_i\cP_i^L\cL_i[\cY_i^{l_1}\cZ_i^{l_2}\cA_i^{l_3+l_4}]|+\omega_i|\E_i\cY_i^{l_1}\cZ_i^{l_2}\cA_i^{l_3+l_4}\cL_i\cP_i^L|+\omega_i|\E_i\nabla(\cP_i^L)^\top\Sigma\nabla(\cY_i^{l_1}\cZ_i^{l_2}\cA_i^{l_3+l_4})|\\
    &=\underbrace{\omega_i|\E_i\cP_i^L\cL_i[\cY_i^{l_1}\cZ_i^{l_2}\cA_i^{l_3+l_4}]|}_{\mathbf{A}}\\
    &+\underbrace{L\omega_i|\E_i\cP_i^{L+1}\cY_i^{l_1}\cZ_i^{l_2}\cA_i^{l_3+l_4+1}|}_\mathbf{B}\\
    &+\underbrace{2L(L+1)\omega_i|\E_i\cP_i^{L+2}\cY_i^{l_1}\cZ_i^{l_2}\cA_i^{l_3+l_4}\x_i^*(N^{-1}A_i)\Sigma(N^{-1}A_i)\x_i|}_{\mathbf{C}}\\
    &+\underbrace{2Ll_1\omega_i|\E_i\cP_i^{L+1}\cY_i^{l_1-1}\cZ_i^{l_2}\cA_i^{l_3+l_4}\x_i^*(N^{-1}A_i)\Sigma(N^{-1/2}B_i\u)|}_{\mathbf{D}}\\
    &+\underbrace{4Ll_2\omega_i|\E_i\cP_i^{L+1}\cY_i^{l_1+1}\cZ_i^{l_2-1}\cA_i^{l_3+l_4}\x_i^*(N^{-1}A_i)\Sigma(N^{-1/2}B_i\u)|}_{\mathbf{E}}\\
    &+\underbrace{4L(l_3+l_4)\omega_i|\E_i\cP_i^{L+1}\cY_i^{l_1}\cZ_i^{l_2}\cA_i^{l_3+l_4-1}\x_i^*(N^{-1}A_i)\Sigma(N^{-1}A_i)\x_i|}_{\mathbf{F}}.
\end{align*}
Here $\mathbf{B}$ and $\mathbf{C}$ arise from evaluating $\cL_i\cP_i^L$, while $\mathbf{D},\mathbf{E}$, and $\mathbf{F}$ arise from computing $\nabla(\cP_i^L)$ and $\nabla(\cY_i^{l_1}\cZ_i^{l_2}\cA_i^{l_3+l_4})$. We begin by bounding $\mathbf{C,D,E}$ and $\mathbf{F}$. By Assumption \ref{ass:A2} and the hypotheses of the lemma, the term $\mathbf{C}$ satisfies
\begin{align*}
    |\mathbf{C}|&\prec\omega_iN^{\delta(L+2)}\Phi^{L+l_2}|\x_i^*(N^{-1}A_i)\Sigma(N^{-1}A_i)\x_i|\\
    &\prec \omega_iN^{\delta(L+2)}\Phi^{L}N^{-1}(N^{-1/2}\|\Sigma^{1/2}A_i\|_F)^2\\
    &\prec N^{-1}(\Phi^2)^{l_4}N^{\delta(L+4)}\Lambda_h\Phi^L\prec c_iN^{\delta(L+4)}\Lambda_h\Phi^L.
\end{align*}
Here, the last step uses $\omega_i\leq |\cC_i|^{l_4}\prec (\Phi^2)^{l_4}\leq 1$.

For $\mathbf{D}$, if $l_1=0$, then $\mathbf{D}=0$, so it remains to consider the case where $l_1\geq 1$. We have
\begin{align*}
    |\mathbf{D}|&\prec\omega_iN^{\delta(L+1)}\Phi^{L+l_2-1}|\x_i^*(N^{-1}A_i)\Sigma(N^{-1/2}B_i\u)|\\
    &\prec\omega_iN^{\delta(L+1)}\Phi^{L-1}N^{-1}\|A_i\|_F\|\Sigma^{1/2}\|_\op N^{-1/2}\|\Sigma^{1/2}B_i\u\|_2\\
    &\prec\omega_iN^{-1/2}N^{\delta(L+2)}\sqrt{\Lambda_h}\Phi^{L}\prec c_iN^{\delta(L+2)}(\Lambda_h+1)\Phi^{L}.
\end{align*}
Here, the last step uses $\omega_i\prec|\u[i]|$ and $\sqrt{\Lambda_h}\leq\Lambda_h+1$. The same arguments show that $|\mathbf{E}|,|\mathbf{F}|\prec c_iN^{\delta(L+4)}\Lambda_h\Phi^L$.

For $\mathbf{B}$, if $l_3+l_4\geq 1$, we have
\[
    |\mathbf{B}|\prec\omega_iN^{\delta(L+1)}\Phi^{L+l_2+1}\prec c_iN^{\delta(L+3)}(\Lambda_h+1)\Phi^L.
\]
Here, we used $\omega_i\prec(|\u[i]|)^{2l_3}|\cC_i|^{l_4}\prec(\u[i]^2+N^{-1})N^{2\delta}(\Lambda_h+1)$. If $l_1\geq 1$, we have
\begin{align*}
    |\mathbf{B}|&\prec\omega_iN^{\delta(L+1)}\Phi^{L+l_2}\E_i|\cY_i|\leq \omega_iN^{\delta(L+1)}\Phi^{L}(\E_i|\cY_i|^2)^{1/2}\\
    &=\omega_iN^{\delta(L+1)}\Phi^{L}(N^{-1}\u^*B_i^*\Sigma B_i\u)^{1/2}\prec c_iN^{\delta(L+2)}(\Lambda_h+1)\Phi^L.
\end{align*}
Here, we used
\[
    (N^{-1}\u^*B_i^*\Sigma B_i\u)^{1/2}\prec N^{-1/2}\|\Sigma^{1/2}B_i\u\|_2\prec N^{-1/2}N^\delta\sqrt{\Lambda_h}\prec N^{-1/2}N^\delta(\Lambda_h+1).
\]
In this case, we also use the fact that $\omega_i\prec|\u[i]|$. Finally, suppose $l_1+l_3+l_4=0$. Then $l_2=L\geq 1$ and $\omega_i\prec 1$, and we have
\[
    |\mathbf{B}|\prec N^{\delta(L+1)}\E_i|\cZ_i|^L|\cA_i|\prec N^{\delta(L+1)}\Phi^{2(L-1)+1}\E_i|\cZ_i|\leq N^{\delta(L+1)}\Phi^L\E_i|\cZ_i|.
\]
Here, we used $2L-1=L+(L-1)\geq L$. We can write
\begin{align*}
    \E_i|\cZ_i|&\leq\E_i|N^{-1/2}\x_i^*B_i\u|^2+|(N^{-1/2}\u^*B_i^\top)\Sigma(N^{-1/2}B_i\u)|\\
    &\prec N^{-1}\|\Sigma^{1/2}B_i\u\|_2^2\prec N^{-1}N^{2\delta}\Lambda_h.
\end{align*}
Combining these estimates shows that $|\mathbf{B}|\prec c_iN^{\delta(L+3)}\Lambda_h\Phi^L$.

For $\mathbf{A}$, the product and chain rules give
\begin{align*}
    &\cL_i[\cY_i^{l_1}\cZ_i^{l_2}\cA_i^{l_3+l_4}]=\cZ_i^{l_2}\cA_i^{l_3+l_4}\cL_i\cY_i^{l_1}+\cY_i^{l_1}\cA_i^{l_3+l_4}\cL_i\cZ_i^{l_2}+\cY_i^{l_1}\cZ_i^{l_2}\cL_i\cA_i^{l_3+l_4}\\
    &+\cA_i^{l_3+l_4}\nabla(\cY_i^{l_1})^\top\Sigma\nabla(\cZ_i^{l_2})+\cZ_i^{l_2}\nabla(\cY_i^{l_1})^\top\Sigma\nabla(\cA_i^{l_3+l_4})+\cY_i^{l_1}\nabla(\cZ_i^{l_2})^\top\Sigma\nabla(\cA_i^{l_3+l_4}).
\end{align*}
Expanding once more with the chain rule and the displayed gradient identities gives
\begin{align*}
    &\mathbf{A}=\underbrace{(-l_1/2-l_2-l_3-l_4)\omega_i\E_i\cP_i^L\cY_i^{l_1}\cZ_i^{l_2}\cA_i^{l_3+l_4}}_{\mathbf{A}_1}\\
    &+\frac12l_1(l_1-1)\omega_i\E_i\cP_i^L\cY_i^{l_1-2}\cZ_i^{l_2}\cA_i^{l_3+l_4}(N^{-1/2}\u^*B_i^\top)\Sigma(N^{-1/2}B_i\u)\\
    &+2l_2(l_2-1)\omega_i\E_i\cP_i^L\cY_i^{l_1+2}\cZ_i^{l_2-2}\cA_i^{l_3+l_4}(N^{-1/2}\u^*B_i^\top)\Sigma(N^{-1/2}B_i\u)\\
    &+2(l_3+l_4)(l_3+l_4-1)\omega_i\E_i\cP_i^L\cY_i^{l_1}\cZ_i^{l_2}\cA_i^{l_3+l_4-2}\x_i^*(N^{-1}A_i)\Sigma(N^{-1}A_i)\x_i\\
    &+2l_1l_2\omega_i\E_i\cP_i^L\cY_i^{l_1}\cZ_i^{l_2-1}\cA_i^{l_3+l_4}(N^{-1/2}\u^*B_i^\top)\Sigma(N^{-1/2}B_i\u)\\
    &+2l_1(l_3+l_4)\omega_i\E_i\cP_i^L\cY_i^{l_1-1}\cZ_i^{l_2}\cA_i^{l_3+l_4-1}\x_i^*(N^{-1}A_i)\Sigma(N^{-1/2}B_i\u)\\
    &+4l_2(l_3+l_4)\omega_i\E_i\cP_i^L\cY_i^{l_1+1}\cZ_i^{l_2-1}\cA_i^{l_3+l_4-1}\x_i^*(N^{-1}A_i)\Sigma(N^{-1/2}B_i\u).
\end{align*}
The arguments for bounding the last six terms are similar to those for $\mathbf{B}-\mathbf{F}$. We therefore omit them and focus on bounding $\mathbf{A}_1$. If $l_3+l_4\geq 1$, then the naive bound gives
\[
    |\mathbf{A}_1|\prec\omega_i N^{\delta L}\Phi^{L+l_2}\prec c_i N^{\delta (L+2)}(\Lambda_h+1)\Phi^L.
\]
Here, as above, we used $\omega_i\prec(\u[i])^{2l_3}|\cC_i|^{l_4}\prec(\u[i]^2+N^{-1})N^{2\delta}(\Lambda_h+1)$. If $l_1\geq 2$, then the naive bound also gives
\[
    |\mathbf{A_1}|\prec\omega_i N^{\delta L}\Phi^{L+l_2}\prec c_i N^{\delta L}(\Lambda_h+1)\Phi^L.
\]
Here, we used $\omega_i\prec|\u[i]|^{l_1}\prec|\u[i]|^2$ and $1\leq \Lambda_h+1$. If $l_2\geq 2$, then we have
\begin{align*}
    |\mathbf{A}_1|\prec\omega_i N^{\delta L}\Phi^{L+l_2-2}\E_i|\cZ_i|\prec c_i N^{\delta(L+2)}(\Lambda_h+1)\Phi^L.
\end{align*}
Here, we used $l_2-2\geq 0$ and, as in the case of $\mathbf{B}$, the bound $\E_i|\cZ_i|\prec N^{-1}N^{2\delta}\Lambda_h$. It remains to consider the case where $l_1,l_2\leq 1$ and $l_3+l_4=0$.

If $l_1=l_2=1,l_3+l_4=0$, then, as in the case of $\mathbf{B}$,
\[
    |\mathbf{A}_1|\prec|\u[i]| N^{2\delta}\Phi^2\E_i|\cY_i|\prec N^{-1/2}|\u[i]|N^{\delta(L+1)}(\Lambda_h+1)\Phi^L.
\]
Here, we used $L=2$. If $l_1=1,l_2+l_3+l_4=0$, then $L=1$, and we can write
\[
    \cP_i=(\beta+N^{-1}\Tr\Sigma A_i)^{-1}-(\beta+N^{-1}\Tr\Sigma A_i)^{-1}\cP_i\cA_i.
\]
Hence,
\begin{align*}
    |\mathbf{A}_1|&\prec|\u[i]||\beta+N^{-1}\Tr\Sigma A_i|^{-1}|\E_i\cY_i|+|\u[i]||\beta+N^{-1}\Tr\Sigma A_i|^{-1}|\E_i\cP_i\cY_i\cA_i|\\
    &\prec |\u[i]|N^{2\delta}\Phi\E_i|\cY_i|\prec N^{-1/2}|\u[i]|N^{3\delta}(\Lambda_h+1)\Phi\leq c_i N^{\delta(L+2)}(\Lambda_h+1)\Phi^L.
\end{align*}
Here, we used $\E_i\cY_i=0$. Similarly, if $l_2=1,l_1+l_3+l_4=0$, then $L=1$, and we have
\begin{align*}
    |\mathbf{A}_1|&\prec |\beta+N^{-1}\Tr\Sigma A_i|^{-1}|\E_i\cZ_i|+|\beta+N^{-1}\Tr\Sigma A_i|^{-1}|\E_i\cP_i\cZ_i\cA_i|\\
    &\prec N^{2\delta}\Phi\E_i|\cZ_i|\prec N^{-1}N^{4\delta}\Lambda_h\Phi\leq c_i N^{\delta(L+3)}\Lambda_h\Phi^L.
\end{align*}
Here, we used $\E_i\cZ_i=0$. Combining all the cases above and using $\Lambda_h\leq\Lambda_h+1$ completes the proof.
\end{proof}

\end{document}
